% Physical unit 011; accepted Session 02.
\GordanUnitStart{011}{printed page 6}
\GordanTopFolio{6}
\[
 a\cos l\,dl\sqrt{1-e^2\cos^2l}=-\sqrt{\sin^2l_0-\sin^2l}\,ds,
\]
\[
 \cos l_0\,dl\sqrt{1-e^2\cos^2l}=\mp\sqrt{\sin^2l_0-\sin^2l}\,\cos l\,d\psi.
\]
In the latter equation the negative or positive sign is to be used according as the longitudes, when the point $\Pi$ recedes from the origin, increase or decrease. On integrating, there results
\[
 s=-a\int_0^l {dl\cos l\sqrt{1-e^2\cos^2l}\over\sqrt{\sin^2l_0-\sin^2l}},
\]
\[
 \psi=\mp\int_0^l {dl\cos l_0\sqrt{1-e^2\cos^2l}\over\cos l\sqrt{\sin^2l_0-\sin^2l}}.
\]

\begin{center}\S\ 7.\end{center}
If the limit $l$ is defined so that it has a value greater than \gq{$-l_0$} and less than \gq{$+l_0$}, then from these formulas there are obtained single appropriate values of both extension and longitude. This agrees with the assertion that the point $\Pi$ in its motion cuts the circle $l$ at only one place; but because the point moves perpetually with constant velocity, this cannot be so. For this reason the formulas proposed by us belong not to the whole extension of the geodetic line $l_0$, but to that part which lies between the origin and the point at which the circle $-l_0$ is touched.

Therefore, in order that they may be transformed into formulas of the general kind, in place of the latitude $l$, which up to now we have taken as variable by itself, there must be introduced a variable of such a kind that it increases according to the recession of the point $\Pi$ from the origin. Such a variable is the angle $\varphi$, defined by the equation
\[
 \sin l_0\cos\varphi=\sin l.
\]

\begin{center}\S\ 8.\end{center}
The angle $\varphi$ rightly contains the qualities that we are seeking. For the relation between the geodetic line and the angle $\varphi$ is this: at the origin $\varphi=0$; at the section with the equator $\varphi={\pi\over2}$; at the tangency with the circle $-l_0$, $\varphi=\pi$; at the second section with the equator $\varphi={3\over2}\pi$; at the second tangency with the circle $l_0$, $\varphi=2\pi$, and so on. Similarly, as the line proceeds, when it touches either circle at those points which we choose to call points of conversion, $\varphi=m\pi$; when it cuts the equator, $\varphi={2m+1\over2}\pi$.

\begin{center}\S\ 9.\end{center}
Since to each point of the geodetic line $l_0$ there correspond certain values of the angle $\varphi$, the quantities
\[
 l\qquad B\qquad T\qquad s\qquad \psi
\]
